\section{Introduction}
\subsection{Identification}
Located at Gokasho, Uji, Kyoto 611-0011, Japan, the Heliotron J project is a research initiative of the Institute of Advanced Energy (IAE) at Kyoto University, which focuses on the advancement of controlled fusion. This experimental facility utilizes a helical-axis heliotron configuration to investigate the magnetic confinement and contribute to understanding how to achieve the stability of high-temperature plasmas. By optimizing magnetic field geometry and advanced heating methods, the project aims to establish the scientific and engineering foundations required for the development of a sustainable, fusion-based power generation system.

This institute was founded in 1996 by the consolidation of the Institute of Atomic Energy and the Plasma Physics Laboratory, but its history traces back to 1941 and its first fusion device ``Heliotron A''. Since then, eight heliotron devices have been built (A, B, C, D, DM, E, DR, and J); the last one, which is the center of this research, was built in 1999 and is now in its 27th year of operation \cite{IAEhistory}. This institute is also host to several other laboratories and research groups, such as the Advanced Energy Generation Division, Advanced Energy Conversion Division, Advanced Energy Utilization Division, Laboratory for Complex Energy Processes and Integrated Research Center for Carbon Negative Science, resulting in a multidisciplinary, multilayered institution \cite{IAEorganization}. Throughout the research stay at IAE, direct project supervision and experimental database access are coordinated by Prof. Shigeru Inagaki, Ph.D. (\texttt{inagaki.shigeru.7s@kyoto-u.ac.jp}, tel. +81-774-38-3484), Professor at the Heliotron J Project. In Costa Rica, this thesis is co-sponsored by the Plasma Laboratory for Fusion Energy and Applications (PlasmaTEC) at the Costa Rica Institute of Technology (\emph{Tecnol\'ogico de Costa Rica}, ITCR). Supervision at PlasmaTEC is provided by thesis advisor Ph.D.~Iv\'an Vargas Blanco alongside committee reader M.Sc.~Ricardo Solano Piedra.

\subsection{Context of the work}

A major obstacle to securing energy via nuclear fusion is magnetohydrodynamic (MHD) instabilities \cite{Adulsiriswad_2024}. In helical devices, MHD oscillations driven by energetic electrons and/or ions lead to significant transport and loss of high-energy particles. These processes reduce heating efficiency and energy confinement, producing serious risk of localized damage to plasma-facing components due to particle hot spots \cite{Zhong2024}. Recent experimental observations in Heliotron J reveal that low-frequency electromagnetic fluctuations are excited during the high-confinement mode (H-mode) phase. These fluctuations couple strongly with edge turbulence, enhancing cross-field heat transport and thereby limiting the expected improvement in energy confinement \cite{Kin2025}.

Such behavior is considered a key factor underlying the H-mode low performance in helical devices compared to tokamaks. In addition, the excitation of large-amplitude MHD oscillations has been observed under certain operational conditions, highlighting the complexity and variability of plasma stability in these systems \cite{Miyashita2025}. Because Heliotron J operates with a flexible $L=1, M=4$ helical-axis configuration assisted by neutral beam injection (NBI) and electron cyclotron heating (ECH), variations in heating schemes and magnetic topologies directly alter the fast-ion stability landscape from discharge to discharge. These findings underscore the urgent need for a comprehensive understanding of both the characteristics and driving mechanisms of MHD oscillations in order to achieve stable plasma operation. Furthermore, traditional manual, shot-by-shot visual inspection of spectrograms is vulnerable to subjective bias and finite-sample artifacts, motivating an automated multi-diagnostic analysis framework. Specifically, this work addresses the research question: Are the dominant coherent oscillations in NBI- and ECH-heated Heliotron~J discharges core-localized, energetic-particle-driven Alfv\'enic modes whose frequency dispersion and nonlinear three-wave coupling respond systematically to electron cyclotron heating? By answering this question, the project provides IAE with a rectified bispectral indexing implementation in \texttt{jpack} and an automated mode-cataloging workflow, alongside transferable multi-probe spectral and higher-order correlation tools for PlasmaTEC.

\subsection{Relevance of the work}

Systematic characterization of MHD oscillations is essential to identify the dominant wave--particle excitation mechanisms and quantify their relationship with plasma heating performance, cross-field transport, and macroscopic energy confinement \cite{Zhong2024, Kin2025, Miyashita2025}. Replacing subjective visual inspection with an automated, statistically calibrated Python framework---combining multi-probe spectral coherence gating, spatial harmonic decomposition, internal radiometric localization, effective-sample-size auditing, and nonlinear bicoherence testing---prevents false-positive detections from background pickup while resolving mode structure and nonlinear dynamics.

The direct beneficiaries of this work are the Heliotron~J experimental team at IAE, Kyoto University, who receive an open-source, modular upgrade to the internal \texttt{jpack} library along with a validated mode catalog template for future experimental campaigns, and PlasmaTEC at ITCR, where this project strengthens national research capabilities in fusion plasma diagnostics, stellarator/heliotron physics, and advanced spectral analysis. The indirect beneficiaries encompass the broader international stellarator and heliotron fusion community investigating energetic-particle-driven MHD stability and anomalous energy transport in three-dimensional magnetic configurations.

\subsection{Scope and Limitations}

This investigation focuses on the spectral, spatial, and nonlinear characterization of energetic-particle-driven modes \cite{Zhong2024}---specifically energetic particle modes (EPMs) and beta-induced Alfv\'en eigenmodes (BAEs)---and quadratic nonlinear three-wave couplings \cite{Kin2025} in Heliotron~J using Mirnov coil, diamagnetic loop, and electron cyclotron emission (ECE) diagnostics across matched $L=1, M=4$ discharges centered on primary shot \#88653 and unpooled replicate shot \#88654; tearing modes and edge-localized modes (ELMs) fall outside scope. Rather than attempting campaign-wide surveys, this project delivers a reusable open-source Python pipeline built atop \texttt{jpack} to establish a reproducible multi-diagnostic workflow.

Spatial and diagnostic constraints bound the resolved mode structures. Because only three toroidal Mirnov probes (MP; MP1, MP3, MP4) are coplanar while MP2 lies off-midplane, toroidal identification is algebraically restricted to modulo~$4$ ($n \equiv -1 \pmod 4$). Poloidally, non-uniform angular spacing of the 14 poloidal Mirnov probes (PMP; PMP1--PMP14) introduces non-uniform discrete Fourier transform (NUDFT) spatial sidelobes requiring phase-concentration validation, and channels PMP1--PMP4 require software polarity inversion for reversed wiring. Moreover, the 16-channel heterodyne ECE radiometer integrates across finite spatial channel widths ($\Delta\rho \approx 0.08\text{--}0.09$), provides uncalibrated relative intensities requiring a low-pass filter (LPF) thermal proxy validated against $W_p/n_e$, and experiences gyrotron radio-frequency (RF) pickup that depresses single-probe edge-to-core coherence during ECH.

Statistical and physical coverage limits further constrain causal inference. Each shot contains only a single ECH switch-off transition ($N=1$ at $291.32\,\mathrm{ms}$) followed by a ${\sim}9\,\mathrm{ms}$ NBI-only phase, restricting frequency resolution and leaving post-ECH bicoherence convergence untestable. Strong serial autocorrelation in smoothed macroscopic profiles ($I_p$, $W_p$, $n_e$) reduces effective sample sizes to $N_{\mathrm{eff}} \approx 4\text{--}47$, rendering empirical correlations suggestive rather than conclusive. Finally, absent local rotational transform ($\iota/2\pi$), calibrated heating powers, Thomson scattering, and divertor Langmuir measurements preclude direct rational-surface resonance or causal transport attribution. Logistically, the three-month duration of the international research stay at Kyoto University and the computational depth of 18-coil by 16-channel multi-diagnostic auditing delimit the scope to detailed validation on the curated \#88653 and \#88654 discharge pair rather than a campaign-wide survey.

\subsection{Objectives}
\subsubsection{General Objective}
To characterize MHD oscillations in the Heliotron J device through Python-based analysis of experimental data, with the purpose of determining their excitation mechanisms and assessing their impact on plasma confinement.

\subsubsection{Specific Objectives}
\begin{itemize}
    \item To determine the spectral properties, spatial structure, and temporal evolution of MHD oscillations in experimental discharges through Python-based signal processing techniques, with the generation of time-frequency representations and the identification of dominant modes.
    
    \item To evaluate the correlation between energetic-particle-driven MHD modes and plasma heating performance through the analysis of synchronization among magnetic fluctuation frequencies, NBI, and ECH, with the derivation of statistical correlation metrics.
    
    \item To investigate nonlinear three-wave coupling between the coherent MHD modes across the heated plasma phases through higher-order bicoherence analysis, with the compilation of statistically validated triad coupling tables and asymptotic convergence profiles.
\end{itemize}